% !TEX root = ../../mixed-sklyanin-cubics.tex
\section{対合を保つ変形と Hochschild 指標の条件付き比較}
\label{mix:sec:equivariant-deformation}

命題\ref{cmo:prop:full-relative-hh}の $c$-族では、対合 $u\mapsto-u$ が
$c$ と $-c$ のファイバーを交換する。
ここでは別のパラメータ $h$ を導入し、対合が各ファイバーを保つ代数族を与える。
中心三次式、滑らかな可換端点、および一般ファイバーの非 PI 性を証明する。
この代数族の相対 smooth proper な残余 dg 模型はまだ構成されておらず、
非可換端点の Hochschild 指標への帰結は、その構成を追加仮定とする。

\subsection{交換子を縮小する族と中心三次式}

$\alpha=(\alpha_1,\alpha_2,\alpha_3)\in\C^3$ に対し、
全生成元の内部次数を一、$h$ の内部次数を零として
\begin{equation}\label{mix:eq:h-relations}
\begin{gathered}
 \mathscr B_\alpha
 =\C[h]\langle x,y,z,u\rangle/
 (r_x,r_y,r_z,[u,x],[u,y],[u,z]),\\
 r_x=[y,z]-h(x^2-\alpha_1u^2),\qquad
 r_y=[z,x]-h(y^2-\alpha_2u^2),\\
 r_z=[x,y]-h(z^2-\alpha_3u^2)
\end{gathered}
\end{equation}
と置く。以下の中心性は、$\alpha_i$ も独立な可換パラメータとした
$\C[h,\alpha_1,\alpha_2,\alpha_3]$ 上で成立する。

\begin{proposition}[交換子パラメータを持つ中心三次式]
\label{mix:prop:h-central}
$\sym(xyz)$ を六つの順列語の平均とすると、
\begin{equation}\label{mix:eq:h-central}
 \Omega_{h,\alpha}
 =\frac{x^3+y^3+z^3}{3}
   -\frac{2h^2}{3}\sym(xyz)
   -u^2(\alpha_1x+\alpha_2y+\alpha_3z)
\end{equation}
は $\mathscr B_\alpha$ の中心元である。
$h=1$ では式\eqref{cmo:eq:mixed-center}に一致する。
\end{proposition}

\begin{proof}
$a=\alpha_1u^2$、$b=\alpha_2u^2$、$c=\alpha_3u^2$ と略記し、
$S=x^3+y^3+z^3$、$T=\sym(xyz)$ と置く。
ここで $a,b,c$ は中心元である。
交換子の積の公式と式\eqref{mix:eq:h-relations}から
\[
 \left[x,\frac{S}{3}-ax-by-cz\right]
 =\frac{h}{3}(yz^2y-zy^2z)
\]
を得る。実際、各積の交換子を展開すれば
\[
\begin{aligned}
 [x,y^3+z^3]&=h(yz^2y-zy^2z-3cy^2+3bz^2),\\
 [x,-by-cz]&=h(cy^2-bz^2)
\end{aligned}
\]
となり、$y^2,z^2$ の補正項が相殺する。
さらに
\[
 yz^2y-zy^2z=[yz,zy]
 =[zy+h(x^2-a),zy]=h[x^2,zy].
\]
六つの順列語を同じ関係式で書き換えると
\begin{equation}\label{mix:eq:h-sym-identity}
 2T=xzy+zyx+h(x^3-ax)
\end{equation}
となる。例えば
\[
\begin{aligned}
 xyz&=xzy+h(x^3-ax),\\
 yxz&=xzy+h(x^3-ax)-h(z^3-cz),\\
 yzx&=zyx+h(x^3-ax),\\
 zxy&=zyx+h(z^3-cz)
\end{aligned}
\]
に $xzy,zyx$ を加えれば式\eqref{mix:eq:h-sym-identity}を得る。
従って
\[
 [x,T]=\frac12[x,xzy+zyx]=\frac12[x^2,zy]
\]
であり、
\[
 [x,\Omega_{h,\alpha}]
 =\frac{h^2}{3}[x^2,zy]
   -\frac{2h^2}{3}\cdot\frac12[x^2,zy]=0.
\]
$x,y,z$ と $a,b,c$ の巡回置換により $y,z$ とも可換である。
$u$ は中心なので、これで全生成元との可換性が従う。
この証明では $h$ による除算を用いていない。
\end{proof}

指定係数
\[
 \alpha^{(1)}=(1,4,9),\qquad \alpha^{(2)}=(16,25,49)
\]
について、二つのブロックの生成元を $(x_j,y_j,z_j)$ ($j=1,2$) とし、
\begin{equation}\label{mix:eq:h-mixed-family}
\begin{gathered}
 \mathscr B=
 \mathscr B_{\alpha^{(1)}}\otimes_{\C[h,u]}
 \mathscr B_{\alpha^{(2)}},\qquad
 w_h=\Omega_{h,\alpha^{(1)}}+100\Omega_{h,\alpha^{(2)}},\\
 \mathscr Q=\mathscr B/(w_h),\qquad
 Q_\lambda=\mathscr Q\otimes_{\C[h]}\C[h]/(h-\lambda)
 \quad(\lambda\in\C)
\end{gathered}
\end{equation}
と置く。$Q_1$ は例\ref{cmo:ex:sixty-four}の $c=0$ の代数 $A_0$ である。
代数の対合
\[
 \iota(h)=h,\qquad \iota(u)=-u,\qquad
 \iota(x_j)=x_j,\quad\iota(y_j)=y_j,\quad\iota(z_j)=z_j
\]
は全ての関係式と $w_h$ を保つので、各 $Q_\lambda$ に作用する。

\subsection{滑らかな Fermat 型の可換端点}

\[
\begin{aligned}
 G_F&=\frac{x_1^3+y_1^3+z_1^3}{3}
       +100\frac{x_2^3+y_2^3+z_2^3}{3},\\
 L_1&=x_1+4y_1+9z_1,\qquad
 L_2=16x_2+25y_2+49z_2,\qquad L_F=L_1+100L_2
\end{aligned}
\]
と置く。$Q_0$ は通常の可換斉次座標環であり、その三次五様体を
\begin{equation}\label{mix:eq:fermat-endpoint}
 Y_F=\operatorname{Proj}Q_0
     =V(G_F-u^2L_F)\subset\mathbb P^6
\end{equation}
と書く。$Y_F$ は、$xyz$ 項を含む
命題\ref{cmo:prop:zero-eckardt-resolution}の指定候補 $X_0$ とは区別する。

\begin{proposition}[可換端点の滑らかさと対合指標]
\label{mix:prop:fermat-eckardt-endpoint}
$Y_F$ は滑らかであり、$p_F=[0:\cdots:0:1]$ は Eckardt 型の点である。
対合 $\iota$ による $\operatorname{Ku}(Y_F)$ の Hochschild homology の
固有空間次元は
\begin{equation}\label{mix:eq:fermat-hh-character}
 \begin{array}{c|rr}
 i&\dim HH_i^+&\dim HH_i^-\\\hline
 -1&5&16\\
 0&2&0\\
 1&5&16
 \end{array}
\end{equation}
であり、その他の次数は零である。
\end{proposition}

\begin{proof}
$F=G_F-u^2L_F$ とする。$u=0$ で六つの active 偏微分が零なら
$x_j=y_j=z_j=0$ ($j=1,2$) となり、射影点を与えない。
$u\ne0$ では $u=1$ と置く。六つの active 偏微分が零なら
\[
 (x_1,y_1,z_1)=(\pm1,\pm2,\pm3),\qquad
 (x_2,y_2,z_2)=(\pm4,\pm5,\pm7).
\]
このとき
\[
 |L_1|=|\pm1\pm8\pm27|\le36,\qquad
 |L_2|=|\pm64\pm125\pm343|
        \in\{154,282,404,532\}.
\]
従って $L_1+100L_2\ne0$ であり、$F_u=-2(L_1+100L_2)\ne0$ となる。
以上により $Y_F$ は滑らかである。

$p_F$ での接超平面は $L_F=0$ であり、その切断は
$V(G_F,L_F)\subset\mathbb P^5$ を底とし $p_F$ を頂点とする cone である。
$Z_F=V(G_F)\subset\mathbb P^5$ は滑らかである。
もし $V(G_F,L_F)$ が特異なら、ある非零の active 座標 $x$ と
$a\in\C$ について $dG_F(x)=a\,dL_F$、$L_F(x)=0$ が成立する。
$Z_F$ の滑らかさから $a\ne0$ であり、$u^2=a$ を選ぶと
$[x:u]$ が $Y_F$ の特異点になる。これは矛盾である。
従って cone の底は滑らかな三次三様体である。

指標には命題\ref{cmo:prop:zero-involution-hh}と同じ同変 residue 計算を用いる。
七変数の二次式の偶・奇固有次元は $(22,6)$ であり、
七つの Jacobian 二次式の固有次元は $(6,1)$ である。
後者の一次独立性は $Y_F$ の滑らかさによる。
従って二次 Jacobian 商の固有次元は $(16,5)$ となる。
residue 同型には $\det(\iota)=-1$ が掛かるので、
$H^{3,2}(Y_F)$ の固有次元は $(5,16)$ である。
複素共役により $H^{2,3}(Y_F)$ も同じであり、
対角 Hodge 成分には恒等に作用する。
HKR と四つの不変な例外線束に対する加法性から
式\eqref{mix:eq:fermat-hh-character}を得る。
\end{proof}

\subsection{一般ファイバーの非 PI 性と有限中心法の制約}

\begin{proposition}[一般の交換子パラメータでは非 PI]
\label{mix:prop:h-generic-nonpi}
$K=\C(h)$ とし、$Q_\eta=\mathscr Q\otimes_{\C[h]}K$ と置く。
この $K$-代数 $Q_\eta$ は非 PI であり、中心上有限ではない。
特に、$h=0,1$ を含む開集合上の族 $\mathscr Q$ に対して、
命題\ref{cmo:prop:full-relative-hh}で用いた
中心上有限な代数という入力は成立しない。
\end{proposition}

\begin{proof}
$u=0$ の一ブロックは
\[
 S_h=K\langle x,y,z\rangle/
 ([y,z]-hx^2,\ [z,x]-hy^2,\ [x,y]-hz^2)
\]
であり、Sklyanin parameters $(a,b,c)=(1,-1,-h)$ を持つ。
\cite[Section~1, pp.~1--2]{GRSklyaninDLL}の曲線と平行移動の規約では
\begin{equation}\label{mix:eq:h-elliptic-data}
 E_h=V(x^3+y^3+z^3+h^2xyz),\qquad
 O=(1:-1:0),\qquad p_h=(1:-1:-h),\qquad
 \sigma_h=t_{p_h}.
\end{equation}
この引用のページ数は同文献の arXiv 版による。
偏微分の計算から $E_h$ は $h^6+27\ne0$ で滑らかであり、
$O$ は flex である。

まず $p_h$ が $K$ 上非捩れであることを示す。
式\eqref{mix:eq:h-elliptic-data}の曲線族は $\C[[h]]$ 上滑らかであり、
その楕円曲線としての単位切断は $O$ である。
$p_h$ は $h=0$ で $O$ に特殊化するが、$O$ の近傍の相対局所座標
$t=z/(x-y)$ では
\[
 t(p_h)=-h/2\ne0.
\]
任意の正整数 $n$ について、この形式群の $n$ 倍写像は
$[n](t)=nt+O(t^2)$ である。
$n$ は $\C[[h]]$ の単元なので $[n]$ は形式群上の同型であり、
$[n]p_h=O$ は $p_h=O$ を強制する。これは表示した座標値に反する。
従って $p_h$ は $\C((h))$ 上、したがって $K$ 上も非捩れである。

次に、$\overline\Omega_h$ を $\Omega_{h,\alpha}$ の $u=0$ への像とする。
これは $\alpha$ に依存しない。
$S=x^3+y^3+z^3$、
$C_+=xyz+yzx+zxy$、$C_-=yxz+xzy+zyx$ と書けば
\[
 C_+-C_-=hS,\qquad
 9\overline\Omega_h=3S-h^2(C_++C_-).
\]
\cite[Section~2, p.~3]{GRSklyaninDLL}における標準的な三次中心元 $g_h=c_3$ の公式へ
$(a,b,c)=(1,-1,-h)$ を代入すると
\begin{align*}
 3g_h
 &=(h^3-1)C_++(h^3+1)C_--2hS\\
 &=h^3(C_++C_-)-3hS=-9h\overline\Omega_h.
\end{align*}
従って
\begin{equation}\label{mix:eq:h-standard-central-scalar}
 g_h=-3h\overline\Omega_h.
\end{equation}
$h\in K^\times$ なので、この比例係数は非零である。

ここで $\mathcal L=\mathcal O_{E_h}(1)$ とする。
平行移動は $\operatorname{Pic}^0(E_h)$ に恒等に作用するため、
$(\sigma_h^*-1)^2[\mathcal L]=0$ であり、
$p_h$ の非捩れ性と $\deg\mathcal L=3$ から
$\sigma_h^*\mathcal L\not\cong\mathcal L$ である。
従ってこれは regular な elliptic triple である。
関係式の双線型化が $\sigma_h$ のグラフを定めるので、
$S_h$ はこの triple に対応する二次代数である。
\cite[Theorem~6.8(i),(ii)]{GRSklyaninATV}の三次核の定理と、
Sklyanin の標準中心元についての
\cite[Section~1, p.~1]{GRSklyaninDLL}の同定により、
$g_h$ は非零の正則な三次中心元であり、
\begin{equation}\label{mix:eq:h-thcr-quotient}
 S_h/(\overline\Omega_h)=S_h/(g_h)
 \cong B(E_h,\mathcal L,\sigma_h).
\end{equation}
ここで右辺は twisted homogeneous coordinate ring である。
式\eqref{mix:eq:h-standard-central-scalar}による除算は一般点上だけで用いており、
$h=0$ の中心性や可換端点の証明には用いていない。

右辺が非 PI であることを、切断の補間によって直接確認する。
標数零では多項式恒等式を多重線型化できるので、
非零の多重線型多項式
\[
 P(X_1,\ldots,X_m)
   =\sum_{\pi\in\mathfrak S_m}
       a_\pi X_{\pi(1)}\cdots X_{\pi(m)}
\]
が恒等的に零にならないことを示せばよい。
変数を並べ替えて $a_{\mathrm{id}}\ne0$ としてよい。
多重線型恒等式は体の拡大で保たれるので、$\overline K$ 上で考える。
$q\in E_h(\overline K)$ を取り、$3d>m$ となる正整数 $d$ を選ぶ。
$\sigma_h$ は無限位数の平行移動なので
\[
 q,\ \sigma_h^d(q),\ \ldots,\ \sigma_h^{(m-1)d}(q)
\]
は相異なる。$\mathcal L_d=
\mathcal L\otimes\sigma_h^*\mathcal L\otimes\cdots\otimes
(\sigma_h^{d-1})^*\mathcal L$ は次数 $3d$ である。
これら $m$ 点の和を $D$ とすると
$H^1(E_h,\mathcal L_d(-D))=0$ だから、
$H^0(E_h,\mathcal L_d)$ から各点でのファイバーへの同時評価は全射である。
従って各 $i$ について、
$\sigma_h^{(i-1)d}(q)$ では非零、残る $m-1$ 点では零となる
$s_i\in H^0(E_h,\mathcal L_d)=B_d$ を選べる。
THCR の積では
\[
 (s_{\pi(1)}\cdots s_{\pi(m)})(q)
 =s_{\pi(1)}(q)\otimes
   s_{\pi(2)}(\sigma_h^d q)\otimes\cdots\otimes
   s_{\pi(m)}(\sigma_h^{(m-1)d}q).
\]
この値が非零となるのは $\pi=\mathrm{id}$ のときだけである。
よって $P(s_1,\ldots,s_m)(q)\ne0$ であり、
$B(E_h,\mathcal L,\sigma_h)$ は非 PI である。

$Q_\eta$ で $u$ と第二ブロックの三変数を零にすると、
$w_h$ は第一ブロックの $\overline\Omega_h$ になる。
従って式\eqref{mix:eq:h-thcr-quotient}への全射
\[
 Q_\eta\twoheadrightarrow
 S_h/(\overline\Omega_h)
 \cong B(E_h,\mathcal L,\sigma_h)
\]
が存在する。PI 性は商へ降りるので $Q_\eta$ も非 PI である。
中心上有限な代数は PI であるため、中心上有限性も否定される。
ある開集合上で $\mathscr Q$ が可換中心部分環上有限なら、
一般点への底変更後も PI となるので、
命題の最後の主張も従う。
\end{proof}

\subsection{相対 smooth proper 性を仮定した指標一定性}

\begin{proposition}[相対模型の存在を仮定した非可換端点の指標]
\label{mix:prop:h-conditional-character}
$R=\C[h,(h^6+27)^{-1}]$ と置く。
$h=0,1$ を含む連結な affine 開集合 $V\subset\operatorname{Spec}R$ 上で、
式\eqref{mix:eq:h-mixed-family}の族に対する残余 dg 模型 $\mathcal C_V$ が
次の性質を満たすと仮定する。
\begin{enumerate}[label=\textnormal{(\roman*)}]
\item $\mathcal C_V$ は $V$ 上相対 smooth proper である。
\item $\mathcal C_V$ は底を固定する $\mu_2$ 作用を持ち、
      二つの端点のファイバーについて
\[
 \mathcal C_0\simeq\operatorname{Ku}(Y_F),\qquad
 \mathcal C_1\simeq\operatorname{Ku}_{\mathrm{coh}}(A_0)
\]
      という $\iota$ と両立する Morita 同値がある。
\end{enumerate}
このとき非可換側 $\operatorname{Ku}_{\mathrm{coh}}(A_0)$ の対合指標も
式\eqref{mix:eq:fermat-hh-character}に一致する。
特に $HH_{\pm1}$ 上の trace はそれぞれ $-11$、
Hochschild homology 全体の交代 trace は $24$ である。
\end{proposition}

\begin{proof}
$V=\operatorname{Spec}R_V$ とする。
Mathew の相対退化定理
\cite[Theorem~1.3]{CMOMathewDegeneration}は、
可換 $\mathbb Q$-代数上の smooth proper dg 圏に適用される。
仮定 (i) により、各 $HH_i(\mathcal C_V/R_V)$ は
有限射影 $R_V$-加群である。
Hochschild 複体の導来底変更とこの射影性から、
各ファイバーは $HH_i(\mathcal C_t)$ に一致する。
仮定 (ii) の作用に対する冪等元
\[
 e_\pm=\frac{1\pm\iota}{2}
\]
は、これらを二つの有限射影直和因子に分ける。
$V$ は連結なので、その階数はそれぞれ一定である。
$h=0$ での階数は命題\ref{mix:prop:fermat-eckardt-endpoint}で計算済みだから、
$h=1$ でも同じである。
trace と交代 trace は、それぞれ $5-16=-11$ と
$2-(-11)-(-11)=24$ から従う。
\end{proof}

\begin{remark}[未解決の相対入力と結論の範囲]
\label{mix:rem:h-relative-gap}
命題\ref{mix:prop:h-conditional-character}の相対模型の存在は、
この節で証明した代数的事実には含まれない。
必要な入力は、$h$-族についての次数別平坦性と相対 Koszul--Gorenstein 計算、
有限線束による proper dg 模型およびその底変更、
さらに $0,1$ を含む開集合上での対角双加群の相対 perfect 性である。
命題\ref{cmo:prop:full-relative-hh}の proper 模型に至る部分を移すことは
これらを調べる具体的な方法だが、その移行をここでは完了したとは扱わない。

両端点の smooth proper 性だけから、同じパラメータ空間上での
相対 smooth proper 性や指標一定性を結論することはできない。
現行の $c$-族では中心上有限な代数を proper な中央スキーム上へ移し、
通常の coherent な対角加群に
命題\ref{cmo:prop:proper-perfect-spread}を適用した。
$h$-族では命題\ref{mix:prop:h-generic-nonpi}によりこの有限中心性が失われる。
従って別の相対対角の構成、またはその perfect 性を保証する入力が必要である。
この障害は $h$-族の相対 smoothness 自体を否定するものではない。

相対模型を構成して指標一致が得られたとしても、それだけでは
$\operatorname{Ku}_{\mathrm{coh}}(A_0)$ と $\operatorname{Ku}(Y_F)$、
あるいは指定候補 $\operatorname{Ku}(X_0)$ との同値は従わない。
その同値には独立に実際の比較核を構成する必要がある。
\end{remark}
